\chapter*{Prólogo}
\addcontentsline{toc}{chapter}{Prólogo}
\markboth{Prólogo}{Prólogo}

En 1977, secuenciar unos pocos miles de bases del bacteriófago $\varphi$X174 fue una hazaña que
llevó años. Hoy, un solo instrumento produce en un día más secuencias de las que existían en todas
las bases de datos del mundo a comienzos de este siglo. La biología se ha convertido, casi sin
pedir permiso, en una ciencia de la información. Quien quiera hacer preguntas sobre genes, células,
poblaciones o ecosistemas necesita hoy dos alfabetos: el de las moléculas y el de los algoritmos.

Este libro nació como acompañante de un curso de bioinformática práctica para estudiantes de
maestría. Cada capítulo corresponde a un módulo del curso, y cada sección a una lección con un
laboratorio ejecutable en Google Colab. Pero el libro está pensado para leerse por sí mismo: allí
donde el notebook muestra \emph{cómo} hacer algo, el texto se detiene a explicar \emph{por qué}
funciona, de dónde viene cada ecuación y qué supuestos la sostienen.

El recorrido tiene una lógica deliberada. Cada tema comienza con una situación cotidiana que
ofrece un punto de apoyo intuitivo; sigue con las definiciones básicas y, paso a paso, sube hasta
la teoría formal que se encuentra en la literatura especializada. Las ecuaciones no aparecen como
adornos: cada una va acompañada de una tabla que explica todos sus símbolos, de un ejemplo
numérico resuelto y de un fragmento de código que la pone a trabajar. Las figuras se generaron a
partir de cálculos reales y reproducibles, y las referencias, que remiten siempre a las fuentes
originales, fueron verificadas una a una.

El libro está organizado en cinco partes. La primera prepara el laboratorio digital y el
vocabulario molecular. La segunda estudia la comparación de secuencias, desde el alineamiento
hasta la filogenética. La tercera sigue a las lecturas de secuenciación desde el instrumento hasta
las variantes. La cuarta aborda la genómica de poblaciones, la transcriptómica, la epigenómica y
la metagenómica. La quinta sube de escala hacia la estructura de las proteínas, las redes, el
aprendizaje automático y los flujos de trabajo reproducibles.

Nada de esto sustituye la práctica. La bioinformática se aprende como se aprende un instrumento:
leyendo la partitura, sí, pero sobre todo tocando. Por eso le invito a leer con el notebook abierto,
a cambiar los parámetros, a romper el código y a repararlo. Si al terminar un capítulo usted puede
explicar a un colega por qué un algoritmo funciona, y no solo cómo se ejecuta, el libro habrá
cumplido su propósito.

\bigskip
\begin{flushright}
{\itshape Juvenal Yosa}\\
{\sffamily\small\color{gris}2026}
\end{flushright}
